\begin{theorem}[Translated vertical six-site geometry]
    \label{thm:peps_vertical_two_plaquette_geometry}
    \lean{TNLean.PEPS.verticalTwoPlaquetteTree_isTree,
        TNLean.PEPS.verticalTwoPlaquetteRegion,
        TNLean.PEPS.verticalTwoPlaquetteIso,
        TNLean.PEPS.verticalTwoPlaquetteTree,
        TNLean.PEPS.verticalTwoPlaquetteTreeDecidableAdj,
        TNLean.PEPS.verticalTwoPlaquetteTree_le,
        TNLean.PEPS.verticalTwoPlaquetteRegion_card,
        TNLean.PEPS.verticalTwoPlaquetteCycleBond,
        TNLean.PEPS.verticalTwoPlaquetteCycleBond_injective,
        TNLean.PEPS.verticalTwoPlaquetteIso_apply,
        TNLean.PEPS.verticalTwoPlaquetteCycleBond_eq_right}
    \leanok
    \uses{thm:peps_translated_two_plaquette_geometry}
    On a finite torus of width at least three and height at least four, the
    translated two-by-three induced region at any vertex $v$ has six distinct
    vertices. It has a spanning path with relative coordinates
    \begin{align}
        (0,0),(0,1),(0,2),(1,2),(1,1),(1,0).
    \end{align}
    Its two remaining bonds are the distinct rightward steps at rows zero and
    one. The assertions include both periodic seams. This is the rotated
    finite-torus geometry for \cite[Theorem~6.16]{Schuch2010PEPS}, with scope
    recorded in \cite{gap:rmp_peps_quantum_double_g_isometry}.
\end{theorem}

\begin{proof}\leanok
    \uses{def:peps_torus_coordinate_swap}
    Coordinate exchange is a graph isomorphism between the two tori. Transport
    the established horizontal induced region, path tree and chords along this
    isomorphism, and evaluate the six vertex coordinates.
\end{proof}

\begin{theorem}[Vertical movement on the original six spins]
    \label{thm:peps_torus_vertical_flux_move}
    \lean{TNLean.PEPS.IsGIsometric.exists_unitary_torusVerticalFluxMove,
        TNLean.PEPS.torusVerticalFluxAssignment,
        TNLean.PEPS.torusVerticalFluxAssignment_eq_treeCycleAssignment,
        TNLean.PEPS.torusVerticalFluxAssignment_right_transport,
        TNLean.PEPS.regularWalkHolonomy_torusVerticalFluxAssignment}
    \leanok
    \uses{thm:peps_vertical_two_plaquette_geometry,
        thm:peps_cycle_assignment_support,
        thm:peps_torus_plaquette_holonomy,
        thm:peps_regular_cycle_global_permutation}
    Retain the preceding period bounds, and let the original torus site tensor
    be $G$-isometric for the regular action of a finite group $G$. Define $u_0(g)$
    to have directed rightward transport $g$ on the lower chord, and define
    $u_1(g)$ to have transport $g$ on both chords. All other bonds are trivial.
    The ordered bond operator is $g^{-1}$ at the horizontal seam. The actual
    counterclockwise plaquette holonomies, ordered from lower to upper, are
    \begin{align}
        (g,1)\quad\longrightarrow\quad(1,g).
    \end{align}
    There is one unitary $W$ on these six original spins, chosen before $g$,
    every boundary configuration $\theta$, and every common exterior and crossing
    assignment $u$, for which
    \begin{align}
        W\Psi_R(u_0(g),\theta)       & =\Psi_R(u_1(g),\theta), \\
        W^*\Psi_R(u_1(g),\theta)     & =\Psi_R(u_0(g),\theta), \\
        (W\otimes I)\Psi(u_0(g),u)   & =\Psi(u_1(g),u),        \\
        (W^*\otimes I)\Psi(u_1(g),u) & =\Psi(u_0(g),u).
    \end{align}
    Thus the adjoint gives downward movement. The global identities concern the
    actual contracted coefficients with internal operators overridden by
    $u_b(g)$; the displayed holonomies concern the literal identity-background
    assignments. This finite-torus realization of
    \cite[Theorem~6.16]{Schuch2010PEPS} has the stated regular-action and period
    restrictions; see \cite{gap:rmp_peps_quantum_double_g_isometry}.
\end{theorem}

\begin{proof}\leanok
    Directed transport cancels the sorting inversion. The four native directed
    steps give the two plaquette holonomies. Recover each literal insertion as
    the derived cycle assignment, and apply the uniform two-cycle operation to
    the actual open and closed contractions. Its global extension is unitary.
    Multiplication by the adjoint gives both reverse identities.
\end{proof}
